\section{结论、评价与适用边界}

\subsection{主要结论}

\begin{enumerate}
\setlength{\itemsep}{2pt}
\setlength{\parsep}{0pt}
  \item 在严格赛前信息边界和整窗标签屏蔽下，Elastic Net 的五折平均 MSE、RMSE 和 MAE 分别为 $220.0855$、$14.8121$ 和 $11.8276$，均优于另外三类候选。$72$ 场需求预测均值为 $170321280$ 人；高关注赛事的不确定性由 OOF 残差分位数单独给出。
  \item 三维 CP-SAT 联合决定 $72$ 场比赛的场馆与时段，得到 $Z_2=0.484946$。独立复算显示无硬约束违反，最小相邻轮休息为 $63$ 小时，场馆负荷为 $2$--$8$ 场；$0.41\%$ 的证书间隙满足预设停止准则，但不构成零间隙最优证明。
  \item 在第二轮结束的共同截面上，联合模拟保持每个样本的晋级名额质量为 $32$。动态方案仅调整 $8/24$ 场，目标由锁定静态方案的 $6.572230$ 提高到 $6.637484$；五种子完整重优化的增益为 $0.065123$--$0.065434$，离散决策一致。
  \item 实际赛程与优化方案存在结构取舍：实际赛程最小休息更长、跨时区和实际旅行更少，优化方案平均休息和黄金覆盖更高。共同边界下的代理综合分为 $0.277058$ 与 $0.480778$，权重和邻居扰动均保持方向；该结果只适用于题内映射，不能解释为 FIFA 现实排程的总体优劣。
\end{enumerate}

\subsection{可核验性}

四问按信息产生顺序连接：问题一只读开球前字段，问题二接收赛前需求，问题三在第二轮结束后更新状态，问题四才引入现实记录和代理。需求预测以“人”为单位进入排程；求解完成后，独立程序从正式 CSV 复算唯一指派、UTC 冲突、休息、容量、安保、日负荷与公平指标。求解器状态说明搜索程度，物理可行性由这些复算结果决定。

静态与动态方案共用动作域、更新环境、价格边界和归一化上下界，因此增益可以直接比较。现实评价则保留数据来源和代理标记：结构量按小时、公里和比例报告，缺失商业字段不伪装成观测。灵敏度分析进一步说明哪些扰动只改变数值，哪些可能改变离散动作，以及哪些结论仍受代理集中限制。

\subsection{局限与使用范围}

需求样本只有 $560$ 条，高观看尾部稀缺，Elastic Net 对极端赛事存在收缩性低估；特征重要性表示预测贡献，不是因果效应。排程模型受题定候选场馆、旅行口径和成本参数约束，未包含真实转播合同、临时管制、训练基地和维护窗口。动态模型采用泊松比分与离散资源档位，并固定第三轮共同截面；实际执行时，更新频率还必须服从通知提前量和资源调度周期。

问题四的商业价值、成本和安保字段来自代理，且映射存在复用集中。真实运营数据补齐前，$Z_4$ 只能描述条件性取舍，不能触发“替换现实排程”的决策。迁移到其他赛事时，可复用的是“需求预测--联合指派--滚动重评--独立审计”的结构；轮次规则、候选场馆、休息下限和资源预算必须按新场景重建。若真实约束无法进入模型或代理稳定性不足，应停止输出总体优劣，改为报告局部方案和待补数据。
